\section{Closedness and strict feasibility of the ramified cone}
\paragraph{An integral bound and closedness.}
For indices starting at zero, put
\[
 (H_3)_{ij}=\frac1{i+j+1}\quad(0\le i,j\le2),\qquad
 (H'_2)_{ij}=\frac1{i+j+2}\quad(0\le i,j\le1),
\]
and $W=\operatorname{diag}(H_3\otimes I_4,H'_2\otimes I_4,
H_3\otimes I_4,H'_2\otimes I_4)$. This is positive definite: its factors
are the moment matrices of $1,s,s^2$ and of $1,s$ with weights $1$ and
$s$ on $[0,1]$. Their inverse traces are $381$ and $54$, respectively,
so $W\succeq I_{40}/381$. The completion terms cancel in the trace, giving
the exact coefficient identity
\begin{equation}
 \left\langle W,\operatorname{diag}(A_G,B_G,A_H,B_H)\right\rangle
 =\int_0^1\operatorname{Tr}(G_J(16s^2)+H_J(16s^2))\,ds
 =8\int_0^1(1+256s^4)\,ds\sum_{k=0}^8c_k
 =\frac{2088}{5}\sum_{k=0}^8c_k.                    \label{eq:ramified-moment}
\end{equation}
Thus $W$ annihilates all $148$ auxiliary coefficient matrices.
If all four Gram blocks are positive semidefinite, bounded source
coordinates bound their total trace and hence every entry. The auxiliary
variables are recovered by fixed linear operations: the $J_k$ are the
upper-right $2\times2$ blocks of the coefficients of $G_J(16s^2)$;
$B_{00},B_{11},B_{01},A_{02}$ are Gram subblocks; and $K_{01},K_{12}$
are the skew parts of $A_{01},A_{12}$. Bounded sources therefore admit
bounded feasible auxiliary witnesses. Taking a subsequence proves that
the projected source cone is closed. For a strictly positive weight
on the whole pencil, use $\widetilde W=\operatorname{diag}(I_9,W)$;
it still annihilates the auxiliaries, and its nine source contractions
are all $2093/5$.

\paragraph{A strict primal point.}
The rational order-$33$ certificate in
\href{https://github.com/anubhav70747/q222-finite-sdp/blob/main/extracted-lifts/canonical-completion/full-polar/slater_proof.txt}{the accompanying strict-feasibility proof} has
$c^0=(1/100,1,1,1,1,1,1,1,1)$ and $J(v)=K(v)$.
For either of its completion polynomials $Q(v)=Q_0+vQ_1+v^2Q_2$,
write its established certificate as
\[
 S=\begin{pmatrix}Q_0&B\\B^T&Q_2\end{pmatrix},\qquad
 T=Q_1-B-B^T=I_4/50.
\]
The two old bounds are $S_G\succeq19I_8/25$ and
$S_H\succeq179I_8/300$. Consequently
\[
 A^0=\begin{pmatrix}
 Q_0&0&16B\\0&16T&0\\16B^T&0&256Q_2
 \end{pmatrix}\succeq\frac8{25}I_{12}.
\]
Set $B^{\rm new}=I_8/100$, and replace the $01$ and $12$ blocks of
$A^0$, and their transposes, by $-I_4/200$. These changes cancel
the additional coefficients $sI_4/100+s^3I_4/100$ from
$sZ_1^TB^{\rm new}Z_1$, preserving the exact polynomial $Q(16s^2)$.
The symmetric perturbation has maximum absolute row sum $1/100$.
Hence $A^{\rm new}\succeq31I_{12}/100$ for both $Q=G,H$.
The full order-$49$ pencil is therefore at least $I_{49}/100$.

In the actual $148$-coordinate parametrization, this point has
\[
 (J_0,J_1,J_2,J_3,J_4)=
 \left(4I_2,0,
       \begin{pmatrix}0&64\\64&0\end{pmatrix},0,
       \begin{pmatrix}0&0\\0&2048\end{pmatrix}\right),
\]
and, on each side,
$B_{00}=B_{11}=I_4/100$, $B_{01}=0$, $A_{02}=16B$,
$K_{01}=K_{12}=0$. The literal auxiliary vector and the four Gram
blocks are supplied in the accompanying rational data.
